\begin{textsample}{NEG Dim 7 – 1990 – Score -1.00 – 1990/tv\_com\_1990\_73.txt}  \label{ex:f7_neg_047}
The commercial opens on a close-up of an older computer monitor.
The screen has a blue background and a centered warning-style label reading “ System Error 86. ” The label is bordered with a diagonal striped pattern, suggesting an alert or system fault.
The scene \textbf{then} cuts to a dimly lit office or workspace.
A person sits at a desk, leaning forward with one hand near their head, as if concentrating or fatigued.
The room contains computer equipment and office furniture.
The lighting is low and dramatic, with blue tones in the background and warmer light around the desk area.
Next, the commercial returns to a close-up of a computer monitor.
Instead of an error message, the screen now shows the face of a man against a blue background.
The image appears like a video image or face displayed on the computer screen.
Only his head and part of his face are visible ; he appears to be looking forward.
The viewpoint \textbf{then} shifts to an \textbf{overhead} angle of the desk setup.
A computer monitor and keyboard are visible on a red-toned desk surface.
The same man ’s face is still visible on the monitor.
A telephone sits to the side.
A person with long hair is shown reclining or leaning back near the desk while holding a telephone handset to their ear.
The framing emphasizes the computer, the phone, and the person ’s position within the workspace.
A close-up follows of a hand resting on the red desk surface.
The hand appears tense or deliberate in placement, with fingers spread slightly.
No object is being actively manipulated in this frame.
Near the end, the image becomes very dark and text appears centered on the screen.
The first word shown is “ There ” in white serif lettering against a mostly black background.
The final frame shows another centered word, “ difference, ” in light-colored text on a dark background.
The sampled frames do not show the full sentence or any product name, and no audio, dialogue, or slogan can be determined from the visuals alone.

% matched lemmas: overhead, then
\end{textsample}
